\documentclass[11pt, letterpaper]{article}

\usepackage{graphicx}
\usepackage{amsmath}
\usepackage[colorlinks=true, allcolors=blue]{hyperref}
\usepackage{fullpage}
\setlength{\parindent}{.25in}
\setlength{\parskip}{8pt}

\title{Capstone Project Literature Review: Analysis of Airline Cost Efficiencies}

\author{Jood Alaswad, Arnol Garcia, Marissa Miller, Brandon Thacker}

\date{\today}

\begin{document}

\maketitle

\section{Introduction}
The United States airline industry operates in a highly competitive and cost sensitive environment. Airlines must continuously balance operating costs, efficiency, and performance while responding to cyclical demand, volatile fuel prices, labor costs, and strong competitive pressure. Airlines strive to minimize expenses while maintaining high levels of capacity utilization and reliability, but the relationship between cost efficiency and operational performance remains unclear. Although low-cost carriers are often perceived as more efficient due to streamlined operations and lower cost structures, legacy carriers can offer advantages in network scope, service differentiation, and connectivity \cite{Barbot2008,Fan2016,GillenLall2004}. This raises an important question: Does cost reduction align with improved operational outcomes, or does it create trade-offs that negatively impact service performance?\\
\indent There is a substantial body of academic literature that examines airline cost structures, business models, and operational performance. Much of this literature focuses on economies of scale, economies of density, operational efficiency, and the strategic differences between low-cost and full-service carriers \cite{Banker1993,GillenLall2004,Lim2015}. Previous research suggests that factors such as load factor, aircraft utilization, network design, labor productivity, and pricing strategy all play significant roles in determining airline efficiency and competitiveness \cite{Barbot2008,Caves1984,Santos2018,Zuidberg2014}. At the same time, other studies have explored service performance metrics such as delays, cancellations, reliability, and broader service quality outcomes in order to highlight the operational challenges airlines face when attempting to maximize efficiency \cite{Lee2023,MayerSinai2002,Yalvac2026}. In hub-and-spoke systems, delays may emerge as a byproduct of network optimization and connection scheduling rather than purely poor performance \cite{MayerSinai2002}. More recent work also shows that airline performance cannot be understood through ticket pricing alone, as ancillary revenues and broader KPI-based cost management now play a major role in airline financial health \cite{Moghadasnian2025,Zhou2020}.\\
\indent Given these considerations, this literature review aims to consolidate existing research on airline cost efficiency and operational performance. There will be a focus on four key areas: (1) the relationship between airline business models and cost efficiency, (2) the role of cost structures and operational efficiency metrics, (3) the relationship between cost minimization and service performance indicators, and (4) how these relationships changed across pre-COVID, COVID, and post-COVID periods \cite{FontanetPerez2022,Gudmundsson2021,Nguyen2022}. By examining these questions, this review provides a foundation for understanding whether cost minimization and service quality are complementary objectives or involve unavoidable trade-offs within the airline industry.

\section{Background}
The modern airline industry has evolved into a complex and competitive market shaped significantly by deregulation, liberalization, and increasing market competition. In the United States, deregulation began in 1978 and intensified pressure on airlines to reduce operating costs while maintaining network coverage and service quality \cite{Borenstein1992}. This shift created greater room for market entry and contributed to the rise of low-cost carriers, which adopted distinct strategies from traditional full-service airlines \cite{Gillen2005,GillenLall2004}.\\
\indent A key feature of the airline industry is the presence of differing business models, most notably between low-cost carriers and legacy airlines, but also hybrid forms that blur the distinction between the two \cite{FontanetPerez2022,Gudmundsson2023}. Low-cost carriers are typically characterized by point-to-point service, high aircraft utilization, dense seating configurations, simplified fleets, aggressive pricing behavior, and reduced service offerings, while legacy carriers usually compete through broader connectivity and more differentiated products \cite{Barbot2008,FontanetPerez2022,GillenLall2004,Gudmundsson2023,Lim2015,OConnell2005}. These differences are not merely marketing choices. Prior literature suggests that business model structure strongly influences how airlines allocate resources, manage capacity, and pursue efficiency.\\
\indent Prior research has shown that cost efficiency is closely tied to traffic density, aircraft utilization, and network design, with evidence suggesting that many airline cost advantages are achieved through improved utilization rather than scale alone \cite{Banker1993,Caves1984,Zuidberg2014}. At the same time, several studies find that low-cost carriers are, on average, more efficient than full-service carriers, though this relationship is influenced by geography, regulation, labor productivity, and firm size \cite{Barbot2008}. More recent work also cautions against treating all airlines as homogeneous units when evaluating efficiency. Airlines differ by region, business model, and operating content, and ignoring these differences can bias results \cite{Li2025}.\\
\indent Operational performance is also a critical component of airline success. Airlines frequently face disruptions, such as delays and cancellations, which affect both reliability and customer satisfaction. However, not all delays should be interpreted as simple signs of inefficiency. Mayer and Sinai (2002) argue that delays at hub airports can reflect rational scheduling and network benefits, as hub carriers cluster flights to maximize connections even when doing so creates congestion \cite{MayerSinai2002}. This suggests that the relationship between cost efficiency and service quality is not always straightforward, and that some operational trade-offs may be inherent to the airline business model itself. Recent research has also emphasized the importance of disruption mitigation, operational continuity, and service quality management, highlighting the trade-offs between minimizing costs and maintaining reliable service \cite{Anupkumar2023,Lee2023,Yalvac2026}.\\
\indent In addition, airline cost and revenue structures have become more complex over time. Airlines increasingly rely on ancillary services such as seat selection and baggage fees as part of broader revenue strategies, meaning that performance can no longer be evaluated solely through ticket fares and passenger volumes \cite{Zhou2020}. As a result, understanding airline efficiency requires attention not only to operating costs, but also to business model design, revenue strategy, and operational performance. These dynamics motivate a more detailed examination of the relationship between cost efficiency and operational outcomes.

\section{COVID-19 and Structural Change}
The COVID-19 pandemic created a major external shock to the airline industry and provides an important context for evaluating changes in cost efficiency and operational performance over time. Prior to the pandemic, airline growth trends were relatively stable, with competition centered largely on cost control, business model strategy, and demand growth. During 2020 and 2021, however, airlines experienced severe demand declines, grounded fleets, operational disruptions, and major financial losses \cite{FontanetPerez2022,Gudmundsson2021}.\\
\indent The literature suggests that the effects of the crisis were not uniform across airline types. Fontanet-P\'erez et al. (2022) find that low-cost and ultra-low-cost carriers in the U.S. generally performed better financially during the first year of the pandemic than full-service network carriers \cite{FontanetPerez2022}. This result suggests that leaner cost structures may provide a short-run resilience advantage under extreme demand shocks, although long-term sustainability remains more uncertain \cite{FontanetPerez2022,Gudmundsson2023}. At the same time, the authors also argue that short-run success under crisis conditions does not necessarily imply that these business models are best adapted to future industry changes, especially in an environment shaped by sustainability pressures and longer-term uncertainty \cite{FontanetPerez2022}.\\
\indent Research on recovery patterns also suggests that COVID should be understood as a severe but likely transitory shock, even if its effects on airline strategy, demand composition, and financial resilience may persist over a longer horizon \cite{FontanetPerez2022,Gudmundsson2021}. Gudmundsson et al. (2021) estimate that passenger demand recovery to pre-COVID levels would take multiple years and vary across regions, but they do not treat the crisis as a permanent break from previous long-run growth patterns \cite{Gudmundsson2021}. This distinction is important for the present study because it supports dividing the analysis into pre-COVID (2016--2019), COVID (2020--2021), and post-COVID (2022--2025) periods. Such a structure makes it possible to examine whether relationships between cost efficiency and operational performance returned to earlier patterns or shifted after the pandemic.\\
\indent More broadly, the COVID literature highlights the importance of resilience in airline business models. It reinforces the idea that cost efficiency, service performance, and strategic fit must be evaluated over time rather than in a single static context. For this reason, COVID is not just a disruption to be acknowledged in passing, but a useful analytical break that helps reveal how different business models respond to extreme stress and recovery.

\section{Methodologies}
The study of airline cost efficiency and operational performance has employed a diverse range of quantitative methodologies, evolving from basic comparative analysis to econometric cost modeling, data envelopment analysis, productivity analysis, and more specialized frontier and network based methods \cite{Banker1993,Barbot2008,Caves1984,Li2025,Wanke2016}. Early research in airline economics often relied on cost function estimation and comparative statistical analysis to distinguish among airlines with different cost structures. Foundational work by Caves et al. (1984) used econometric methods to separate economies of density from economies of scale, establishing an important framework for understanding why airline costs differ across business types and route systems \cite{Caves1984}.\\
\indent Subsequent research expanded on these approaches by incorporating more detailed operational variables and firm level data. Banker and Johnston (1993) applied empirical cost modeling to identify major cost drivers in the U.S. airline industry, emphasizing the role of output, operational structure, and network design in explaining variation in costs \cite{Banker1993}. Barbot et al. (2008) combined DEA and total factor productivity methods to compare airlines across business models and regions, showing that methodological choice can influence how productivity and efficiency differences are interpreted \cite{Barbot2008}. More recent studies, such as Zuidberg (2014), have used regression based econometric approaches to analyze airline cost economies while incorporating variables such as load factor, fleet utilization, and aircraft characteristics \cite{Zuidberg2014}. These methods allow researchers to isolate the effects of individual operational variables while controlling for differences across carriers.\\
\indent In addition to traditional econometric approaches, modern research has increasingly incorporated efficiency measurement techniques such as data envelopment analysis (DEA). Nguyen et al. (2022) use a two stage DEA approach to compare revenue efficiency across airline business models. Their study combines a non-convex meta-frontier with truncated regression and bootstrapping to avoid some of the bias associated with standard DEA methods \cite{Nguyen2022,SIMAR1998,SIMAR2000}. Kaya et al. (2023) employ a super efficiency slack based measure approach to evaluate multiple dimensions of airline efficiency, including customer, financial, internal process, learning, and environmental measures \cite{KAYA2023}. Li et al. (2025) further argue that airline heterogeneity must be incorporated into efficiency analysis rather than assuming all airlines operate under the same production frontier \cite{Li2025}. Related work in transport efficiency has also used frontier based methods under uncertainty, such as Fuzzy-DEA, to improve the robustness of performance measurement \cite{Wanke2016}. These approaches are valuable because they broaden efficiency measurement beyond simple accounting ratios and allow for richer comparisons across firms.\\
\indent More recent research has also highlighted the importance of heterogeneity in airline efficiency analysis. Li et al. (2025) argue that airlines should not be treated as homogeneous decision making units because they differ across region, business model, and operational focus. Using a meta-frontier DEA approach, they divide airline performance into operation, service, and sales stages and find that much of airline inefficiency originates in the operations stage \cite{Li2025}. This finding is particularly relevant for the present study because it suggests that operational cost and operational performance may be more directly linked than broader financial measures alone imply.\\
\indent Research on operational reliability and service performance has also adopted quantitative modeling approaches. Mayer and Sinai (2002) use large scale U.S. flight data to examine delays and conclude that congestion at hubs can emerge as an equilibrium outcome of network optimization rather than purely market failure \cite{MayerSinai2002}. Lee (2023) applies analytical modeling to evaluate how airlines invest in mitigating operational disruptions, highlighting the trade-offs between reducing losses and maintaining reliability \cite{Lee2023}. Together, these studies suggest that service quality metrics such as delays and cancellations are useful, but must be interpreted in relation to network structure and operational strategy.\\
\indent Despite the variety of methodologies employed, several limitations persist. Many studies rely on observational data, which may introduce bias or limit causal interpretation. In addition, differences in variable definitions, business model classification, and modeling choices can make results difficult to compare across studies. In the present study, these ideas are applied using BTS datasets that capture airline financial activity, traffic, and on-time performance. Rather than estimating a full production frontier, the analysis focuses on constructing comparable performance indicators from BTS Form 41, T-100, and on-time data to evaluate how operating cost patterns relate to load factor, passenger volume, delays, cancellations, and reliability across airline business models \cite{bts_ontime, bts_transtats}. This approach aligns the project more closely with the strand of literature that links efficiency to observable operational outcomes rather than treating airlines only as abstract production units. 

\section{Conclusion}
The airline industry presents a complex environment in which cost efficiency and operational performance are closely intertwined, but not always aligned. The existing literature demonstrates that while airlines can achieve efficiency gains through improved utilization, network design, pricing strategy, and cost management, these strategies can also introduce challenges related to service reliability, operational disruptions, and customer facing service quality \cite{Anupkumar2023,Barbot2008,Caves1984,GillenLall2004,Lee2023,MayerSinai2002,Zuidberg2014}. In addition, the literature suggests that airline performance differs significantly across business models, and that low-cost carriers are often more efficient than full-service carriers, though this relationship is shaped by heterogeneity in geography, scale, and operational design \cite{Barbot2008,Li2025}.
\indent More recent research has further shown that airline cost efficiency cannot be understood through ticket pricing alone. Revenue strategies, ancillary services, and multi-dimensional KPI frameworks all affect airline financial health and complicate the relationship between costs and performance \cite{Moghadasnian2025,Zhou2020}. At the same time, the COVID-19 pandemic revealed that business models respond differently to severe external shocks, making it especially important to examine cost-performance relationships across pre-COVID, COVID, and post-COVID periods \cite{FontanetPerez2022,Gudmundsson2021}.\\
\indent By incorporating previous research on business models, cost structures, efficiency measurement, service quality, and post-COVID resilience, this study establishes a broader foundation for examining these relationships using empirical data \cite{FontanetPerez2022,Gudmundsson2021,KAYA2023,Li2025,Nguyen2022,Yalvac2026}. The findings from this analysis aim to contribute to a better understanding of whether cost reduction strategies enhance overall performance or create unintended consequences, ultimately informing how airlines can more effectively balance efficiency and service quality in a competitive market.

\bibliographystyle{abbrv}
\bibliography{refs}

\end{document}